\section*{अध्याय \ref{ch:b3:developmental-genetics} --- परिवर्धनी आनुवंशिकी और आकार का विकास}

\begin{solution}{exo:b3:developmental-genetics:1}
मातृ-प्रभाव: \emph{bicoid}, जिसमें उत्परिवर्ती माँओं के भ्रूणों में सिर
तथा वक्ष नहीं होते। रिक्ति: \emph{Krüppel}, जिसमें खंडों का कोई सटा
हुआ गुच्छा ग़ायब। युग्म-नियम: \emph{fushi tarazu} (या
\emph{even-skipped}), जिसमें हर दूसरा खंड ग़ायब। खंड-ध्रुवता:
\emph{hedgehog} (या \emph{engrailed}), जिसमें हर खंड का एक हिस्सा बाकी
के किसी दर्पण-प्रतिबिंब से बदल जाता है।
\end{solution}

\begin{solution}{exo:b3:developmental-genetics:2}
कोई ऐसा जीन जिसका उत्पाद माँ अंडे में जमा कर देती है, इसलिए उस भ्रूण
का लक्षणरूप उसके जीनरूप के पीछे चलता है। उसके भ्रूणों में सिर तथा वक्ष
नहीं होते और दोनों सिरों पर पश्च संरचनाएँ होती हैं; वे मर जाते हैं। वह
पैतृक वन्य-प्ररूपी युग्मविकल्पी युग्मनज में है, पर \emph{bicoid} का
mRNA अंडजनन के दौरान माँ की धात्री कोशिकाएँ बनाती हैं और वह निषेचन से
पहले ही उस ध्रुव पर बाँध दिया जाता है; युग्मनज की अपनी प्रति घंटों बाद
अनुलेखित होती, यानी उस प्रवणता का समय बीत जाने के बाद।
\end{solution}

\begin{solution}{exo:b3:developmental-genetics:3}
\omterm{def:b3:developmental-genetics:hox}{सहरैखिकता}: गुणसूत्र पर \omterm{def:b3:developmental-genetics:hox}{Hox जीनों} का क्रम शरीर पर उनके
अभिव्यक्ति-प्रांतों के क्रम से मेल खाता है। पश्च प्रबलता: जहाँ कई
अभिव्यक्त हों वहाँ सबसे पश्च वाला पहचान तय करता है।
\emph{Hoxc8}-विहीन कोई चूहा कोई अग्र रूपांतरण दिखाता है --- उसकी पहली
कटि-कशेरुका किसी वक्षीय की पहचान ले लेती है और कोई पसली ढोती है। सिर
में \emph{Antennapedia} स्पर्शश्रृंगों को टाँगों में बदल देता है, यानी
दूसरे वक्षीय खंड की पहचान में।
\end{solution}

\begin{solution}{exo:b3:developmental-genetics:4}
किसी कलमित पृष्ठीय ओष्ठ ने कोई दूसरा पूरा अक्ष प्रेरित किया, जो
अधिकांशतः परपोषी कोशिकाओं का बना था: यानी उस कलम ने अपने पड़ोसियों को
निर्देश दिया --- कोई \omterm{def:b3:developmental-genetics:organiser}{संगठक}। उसके अणु स्रवित BMP-विरोधी हैं (कॉर्डिन,
नॉगिन, फ़ॉलिस्टैटिन)। बाह्यचर्म तब तंत्रिकीय बनता है जब BMP-संकेतन
अनुपस्थित हो और तब बाह्यत्वचा जब BMP मौजूद हो, इसलिए तंत्रिकीय वही है
जो बाह्यचर्म तब करता है जब उसे कुछ बताया ही न जाए: वह \omterm{def:b3:developmental-genetics:organiser}{संगठक} किसी निर्देश
को चुप कराकर काम करता है।
\end{solution}

\begin{solution}{exo:b3:developmental-genetics:5}
$k = \ln 2/(35\times 60) = \qty{3.3e-4}{s^{-1}}$; $\lambda = \sqrt{4/
3.3\times 10^{-4}} = \qty{110}{\micro m}$; अनुपात $\mathrm{e}^{250/110}
= \mathrm{e}^{2.27} = 9.7$।
\end{solution}

\begin{solution}{exo:b3:developmental-genetics:6}
चार प्रतियाँ $c_{0}$ दुगुना कर देती हैं: वह सीमा $\lambda\ln 2 =
\qty{69}{\micro m}$ पीछे खिसककर \qty{309}{\micro m} पर आ जाती है; एक
प्रति उसे आधा कर देती है: \qty{69}{\micro m} आगे, \qty{171}{\micro m}
पर। हर खिसकाव अंडे की लंबाई का \qty{14}{\%} है --- जो देखे गए से बड़ा
है, और वही उन तंत्रों का प्रमाण है जो उस मात्रा-निर्भरता को दबा देते
हैं।
\end{solution}

\begin{solution}{exo:b3:developmental-genetics:7}
$\delta x = \lambda\,\delta c/c$: \qty{8}{\micro m} के लिए $\delta c/c =
8/100 = \qty{8}{\%}$। किसी प्वासों गिनती $N$ की सापेक्ष
त्रुटि $1/\sqrt{N}$ है, इसलिए $N = 1/0.08^{2} \approx 160$ अणु गिनने
पड़ेंगे --- यानी किसी केंद्रक में मौजूद अणुओं का कुछ ही प्रतिशत, या
समय भर समाकलित कोई एक गिनती।
\end{solution}

\begin{solution}{exo:b3:developmental-genetics:8}
$k = \ln 2/35 = \qty{0.0198}{min^{-1}}$: $1 - \mathrm{e}^{-kt}$ देता है
\qty{60}{min} पर $0.70$, \qty{90}{min} पर $0.83$, \qty{120}{min} पर
$0.91$, और \qty{150}{min} पर $0.95$। वह प्रवणता पढ़े जाते समय भी कुछ
प्रतिशत चढ़ ही रही होती है, जो उस पठन-खिड़की भर में किसी सीमा को
$\lambda\ln(1/0.9) \approx \qty{10}{\micro m}$, यानी लगभग एक केंद्रक,
खिसका देता: उस भ्रूण को या तो किसी नियत समय पर पढ़ना पड़ता है या कोई
ऐसा पठन काम में लेना पड़ता है जो कुल स्तर के प्रति असंवेदी हो।
\end{solution}

\begin{solution}{exo:b3:developmental-genetics:9}
कोई दूसरा ZPA कोई दर्पण-प्रतिबिंबी दुहराव देता है, 4-3-2-2-3-4: अब वे
अग्र कोशिकाएँ ऊँचा Shh देखती हैं और पश्च अंगुलियाँ बन जाती हैं। थोड़ा
Shh छोड़ता कोई मनका कोई आंशिक दुहराव देता है --- कोई अतिरिक्त अग्र
अंगुली 2, मान लीजिए 2-2-3-4 --- क्योंकि वे अग्र कोशिकाएँ केवल वह कम
सांद्रता देखती हैं जो अंगुली 2 निर्दिष्ट करती है।
\end{solution}

\begin{solution}{exo:b3:developmental-genetics:10}
विशिष्ट हल $s/k$; समघात हल $\cosh(x/\lambda)$,
$\sinh(x/\lambda)$; $0$ पर कोई फ्लक्स न होना $\sinh$ को मार देता है;
और $c(L) = 0$ वह आयाम तय कर देता है:
$c(x) = (s/k)\bigl[1 - \cosh(x/\lambda)/\cosh(L/\lambda)\bigr]$।
वह परिच्छेदिका $x = 0$ के पास कोई पठार है जो उस गर्त पर $\lambda$
चौड़ाई की किसी सीमा-परत में शून्य तक गिरता है --- यानी किसी बिंदु से
आता कोई चरघातांक नहीं, बल्कि किनारे वाला कोई समतल क्षेत्र; और सबसे ऊँची
सांद्रता उस गर्त से सबसे दूर है। $L \ll \lambda$ के लिए $\cosh u \approx 1 +
u^{2}/2$ और $c \approx s(L^{2} - x^{2})/(2D)$: यानी $k$ से स्वतंत्र
कोई परवलय, क्योंकि अणु अपघटित होने से पहले ही उस गर्त तक पहुँच जाते हैं।
\end{solution}

\begin{solution}{exo:b3:developmental-genetics:11}
$x_{1} = \lambda\ln(1/0.3) = 1.20\lambda$ और $x_{2} = \lambda\ln(1/0.08)
= 2.53\lambda$: पहला क्षेत्र $1.20\lambda$ चौड़ा है, दूसरा
$1.32\lambda$, और दोनों $c_{0}$ से स्वतंत्र, जो उन्हें साथ खिसका देता
है। तीसरा क्षेत्र, $L - 2.53\lambda$, अंडे की लंबाई का सारा फेर सोख
लेता है: $\lambda = \qty{100}{\micro m}$ के साथ \qty{450}{\micro m} के
किसी अंडे में वह उदर \qty{197}{\micro m} है, और \qty{550}{\micro m} के
किसी अंडे में \qty{297}{\micro m} --- यानी वह प्रतिरूप मापक्रमित नहीं
है। तंत्र: पश्च ध्रुव से आती कोई दूसरी प्रवणता (Nanos, Caudal, या वह
सिरा-तंत्र), जो Bicoid के साथ पढ़ी जाए, ताकि स्थितियाँ उस अनुपात से तय
हों; या कोई ``विस्तारक'' अणु, जो वहाँ बने जहाँ वह आकारजन कम हो और
$\lambda$ को तब तक लंबा करे जब तक वह प्रवणता उस अंडे को भर न दे।
\end{solution}

\begin{solution}{exo:b3:developmental-genetics:12}
कोई कूटलेखी बदलाव उस प्रोटीन को हर उस जगह बदल देता है जहाँ वह काम करता
है: \emph{Pitx1} \omterm{def:b3:endocrinology:axes}{पीयूष} तथा जबड़ा भी बनाता है, और \emph{sonic hedgehog}
तंत्रिका नलिका, आँत तथा दाँतों का प्रतिरूपण भी करता है, इसलिए कोई टूटा
प्रोटीन घातक है या बहुप्रभावी ढंग से अपंग करने वाला। कोई संवर्धक किसी
एक ऊतक में किसी एक समय अभिव्यक्ति चलाता है; उसे विलोपित करने से वह
संरचना हट जाती है जो वह ऊतक बनाता, और उस प्रोटीन का हर दूसरा उपयोग अछूता
रह जाता है। संवर्धक मॉड्युलर होते हैं और प्रायः आंशिक रूप से
अनावश्यक-बहुल, इसलिए बीच के चरण जीवनक्षम हैं, और वही रास्ता --- कोई
स्विच खोना --- स्टिकलबैक की अनेक झीलों में स्वतंत्र रूप से लिया गया है।
\end{solution}

\begin{solution}{pb:b3:developmental-genetics:1}
\textbf{1.} $k = \ln 2/(35\times 60) = \qty{3.3e-4}{s^{-1}}$; $1/k =
\qty{3030}{s} = \qty{50}{min}$।
\textbf{2.} $\lambda = \sqrt{4/3.3\times 10^{-4}} = \qty{110}{\micro m}$;
वह अंडा $4.5\lambda$ का है।
\textbf{3.} ध्रुव से ध्रुव तक $\mathrm{e}^{4.5} \approx 90$;
$\mathrm{e}^{8.5/110} = 1.08$, यानी प्रति केंद्रक \qty{8}{\%} का बदलाव।
\textbf{4.} $J = 50\times\sqrt{4\times 3.3\times 10^{-4}} = 50\times
0.036 = \qty{1.8}{nM\,\micro m\,s^{-1}}$। एक nM प्रति
\unit{\micro m^3} $0.6$ अणु है, इसलिए प्रति \unit{\micro m^2}
प्रति सेकंड लगभग $1.1$ अणु।
\textbf{5.} केंद्रक का आयतन $\tfrac{4}{3}\pi\times 27 = \qty{113}{\micro
m^3}$। \qty{50}{nM} पर प्रति \unit{\micro m^3} $30$ अणु:
यानी \num{3400} अणु। अंडे के बीच $c = 50\,\mathrm{e}^{-2.27} =
\qty{5.2}{nM}$: यानी लगभग \num{350} अणु।
\textbf{6.} \qty{60}{min} पर $kt = 1.19$: $0.70$; \qty{120}{min} पर
$kt = 2.38$: $0.91$। पढ़े जाते समय वह प्रवणता स्थायी के \qty{10}{\%}
के भीतर है --- फिर भी इतनी चढ़ रही है कि किसी सीमा को
$\lambda\ln(1/0.91) \approx \qty{10}{\micro m}$, यानी एक केंद्रक,
खिसका दे।
\textbf{7.} $x_{1} = 110\ln(1/0.3) = \qty{132}{\micro m}$
(\qty{26}{\%}); $x_{2} = 110\ln(1/0.08) = \qty{278}{\micro m}$
(\qty{56}{\%})।
\textbf{8.} चार प्रतियाँ: $c_{0} = \qty{100}{nM}$; दोनों सीमाएँ
$\lambda\ln 2 = \qty{76}{\micro m}$ पीछे खिसककर $208$ और
\qty{354}{\micro m} पर आ जाती हैं। छह प्रतियाँ: $\lambda\ln 3 = \qty{121}{\micro m}$,
यानी $253$ और \qty{399}{\micro m} पर।
\textbf{9.} एक प्रति: \qty{76}{\micro m} आगे, यानी $56$ और
\qty{202}{\micro m} पर। सिर-क्षेत्र $132$ से सिकुड़कर
\qty{56}{\micro m} रह गया; वक्ष ने अपनी चौड़ाई (\qty{146}{\micro m})
रखी और आगे सरक गया; और उदर $222$ से बढ़कर \qty{298}{\micro m} हो गया।
\textbf{10.} $\delta x = 110\times 0.1 = \qty{11}{\micro m}$, यानी $1.3$
केंद्रक-अंतराल।
\textbf{11.} एक केंद्रक: $\delta c/c = 8.5/110 = \qty{7.7}{\%}$।
\qty{10}{\%} से, $n = (10/7.7)^{2} \approx 1.7$: यानी दो स्वतंत्र पठन
(या दो पड़ोसी केंद्रकों का औसत) काफ़ी हैं।
\textbf{12.} $1/\sqrt{350} = \qty{5.3}{\%}$: किसी एक तात्क्षणिक गिनती
की परिशुद्धि पहले ही मानी गई \qty{10}{\%} से बेहतर है, इसलिए देखा गया
शोर अकेली गिनती का शोर नहीं बल्कि बंधन-बलगतिकी तथा उस पठन-तंत्र से आता
है --- और समय भर औसत लेने से वह और सुधरता है।
\textbf{13.} अंडे की लंबाई का $278/450 = \qty{62}{\%}$ और
$278/550 = \qty{51}{\%}$: यानी \qty{11}{\%} का अंतर, छह केंद्रक। नियत
$\lambda$ के साथ वह प्रतिरूप मापक्रमित नहीं होता।
\textbf{14.} यदि $D \propto L^{2}$ हो, जिसमें $k$ नियत हो, तो
$\lambda \propto L$ और $x_{T}/L = (\lambda/L)\ln(c_{0}/T)$ हर अंडे में
वही है: यानी पूर्ण मापक्रमण। पर $D$ उस अणु और उस कोशिकाद्रव्य का गुण
है, उस अंडे के आकार का नहीं, इसलिए जैसा कहा गया वैसा यह अप्रशंसनीय है;
और असली भ्रूणों में मापक्रमण आंशिक है तथा दूसरे पठनों से आता है।
\textbf{15.} $c^{5}/(c^{5} + T^{5}) = 0.1$ तब जब $c/T = 9^{-1/5} = 0.64$
और $0.9$ तब जब $c/T = 9^{1/5} = 1.55$: यानी सांद्रता में $2.4$ का गुणक,
$\Delta x = 110\ln 2.4 = \qty{97}{\micro m}$ --- यानी ग्यारह केंद्रक,
जो देखी गई दो-तीन केंद्रक की सीमा से कहीं चौड़ा है; अकेला उस प्रवणता का
पठन इतना तीखा नहीं है।
\textbf{16.} तीखापन किसी दी सांद्रता को अधिक निर्णायक ढंग से चालू या
बंद पर मानचित्रित करता है, इसलिए वह संक्रमण-क्षेत्र सँकरा होता है; पर
सांद्रता में \qty{10}{\%} की त्रुटि उस बिंदु को, जहाँ वह देहली पार होती
है, फिर भी $\lambda\,\delta c/c$ खिसका देती है: वह अनुक्रिया कोई पूर्ण
सोपान हो सकती है और वह सोपान फिर भी ग़लत जगह पड़ सकता है।
\textbf{17.} यदि दोनों चालू होते, तो हर एक दूसरे को दबा रहा होता,
इसलिए कम से कम एक बंद कर दिया जाता; सहकारी दमन के साथ केवल
एक-चालू-एक-बंद ही स्थिर अवस्थाएँ हैं, और उनके बीच की सीमा कोई स्विच है,
कोई ढलान नहीं।
\textbf{18.} $\mathrm{d}\ln\lambda = \tfrac{1}{2}(0.02 - 0.06) = -0.02$
प्रति अंश: $\lambda$ प्रति अंश \qty{2}{\%} गिरता है, यानी
\qty{5}{\celsius} के लिए \qty{10}{\%}, और \qty{99}{\micro m} पर आ जाता
है; और \emph{hunchback} की सीमा $278$ से \qty{250}{\micro m} पर, यानी
तीन केंद्रक आगे, खिसक जाती है। मक्खियाँ \qty{18}{} और
\qty{29}{\celsius} के बीच सामान्य अनुपातों के साथ विकसित होती हैं,
इसलिए कुछ न कुछ भरपाई करता है --- कोई ऐसा स्रोत या पठन जो तापमान के
साथ दूसरी ओर खिसके।
\textbf{19.} $2^{8} = 256$ कूट; चौदह काम में। पश्च प्रबलता के साथ केवल
सबसे पश्च चालू जीन मायने रखता है, इसलिए अधिक से अधिक नौ भिन्न पहचानें
(आठ जीन, या कोई नहीं): वह कूट कोई क्रमांकन है, कोई द्विआधारी शब्द नहीं।
\textbf{20.} तीसरा वक्षीय खंड दूसरा बन जाता है: \emph{Ubx} उस
पंख-कार्यक्रम को दबाए था जो कोई वक्षीय खंड चूक से चलाता है, और उस
दमनकारी को हटाने से वह पुरखा अवस्था छूट जाती है --- यानी चार पंख, जैसे
उस मक्खी के चार-पंखी पूर्वजों के थे।
\textbf{21.} हंस का ग्रैव क्षेत्र कशेरुका 22 तक चलता है: यानी अनेक
कशेरुकाओं की कोई लंबी गर्दन, जबकि चूहे की सात। वह जीन वही है; उसकी
अभिव्यक्ति की अग्र सीमा पश्च की ओर खिसक गई --- यानी इस बात का कोई
नियामक बदलाव कि कोई \omterm{def:b3:developmental-genetics:hox}{Hox जीन} कहाँ चालू होता है।
\textbf{22.} अब हर किनारा कोई स्रोत है; Shh दोनों किनारों पर सबसे ऊँचा
है (अंगुली 4), बीच की ओर गिरता है (3, फिर 2), और उस कम स्तर तक कभी
नहीं पहुँचता जो अंगुली 1 निर्दिष्ट करता है, इसलिए बीच में पीठ से पीठ
जोड़े दो अंगुलियाँ 2 पड़ती हैं: 4-3-2-2-3-4।
\textbf{23.} \emph{Pax6} कोई संरक्षित स्वामी नियामक है जो वह भी
नेत्र-कार्यक्रम छेड़ सकता है जो उस परपोषी के पास हो: उस मक्खी ने कोई
मक्खी-आँख बनाई, इसलिए वह चूहा-जीन कोई चूहा-आँख की सूचना नहीं ढोता, केवल
यह आदेश ढोता है कि ``यहाँ कोई आँख बनाओ''। यह उस जीन-जाल की \omterm{def:b3:developmental-genetics:evodevo}{गहन समजातता}
दिखाता है, उन अंगों की समजातता नहीं: संयुक्त आँख और कैमरा आँख किसी ऐसे
पूर्वज से अलग-अलग विकसित हुईं जिसके पास \emph{Pax6} वाला कोई
प्रकाश-ग्राही टुकड़ा था।
\textbf{24.} \qty{2}{mm} भर पाँच अंगुलियाँ \qty{0.4}{mm} का कोई
तरंगदैर्घ्य हैं; छह के लिए \qty{0.33}{mm} चाहिए, यानी छठे भाग की कमी।
चूँकि $\lambda \propto \sqrt{D_{\text{निरोधी}}/k_{\text{निरोधी}}}$, उस
निरोधक का क्षय \qty{44}{\%} बढ़ाइए या उसका विसरण \qty{31}{\%} घटाइए;
और उस हाथ-पट्टिका में पश्च \omterm{def:b3:developmental-genetics:hox}{Hox जीनों} की मात्रा घटाने से वही होता है, और
उन उत्परिवर्ती चूहों की छह, सात या उससे अधिक पतली अंगुलियाँ होती हैं।
\textbf{25.} $\lambda \approx \qty{110}{\micro m}$; स्रोत का दुगुना
होना हर सीमा को \qty{76}{\micro m} खिसका देता है; और सांद्रता में
\qty{10}{\%} शोर \qty{11}{\micro m}, यानी लगभग एक केंद्रक, की स्थानिक
त्रुटि है।
\end{solution}
